\documentclass[11pt]{amsart}
\usepackage{amsmath,amssymb,amsthm,booktabs}
\usepackage[margin=1.15in]{geometry}
\usepackage[hidelinks]{hyperref}
\newtheorem{theorem}{Theorem}[section]
\newtheorem{lemma}[theorem]{Lemma}
\newtheorem{corollary}[theorem]{Corollary}
\theoremstyle{remark}
\newtheorem{remark}[theorem]{Remark}
\newcommand{\Z}{\mathbb{Z}}
\newcommand{\PP}{\mathbb{P}}
\newcommand{\E}{\mathbb{E}}
\newcommand{\I}{\mathbb{I}}
\newcommand{\doi}[1]{\href{https://doi.org/#1}{doi:\nolinkurl{#1}}}
\newcommand{\arxiv}[1]{\href{https://arxiv.org/abs/#1}{arXiv:\nolinkurl{#1}}}
\setlength{\emergencystretch}{2.5em}

\title[Two-sided conditioning in random walk in random scenery]{Good configurations of random walk in random scenery
under two-sided conditioning}
\author{Yitzchak Shmalo}
\address{Einstein Institute of Mathematics, The Hebrew University of Jerusalem, Jerusalem, Israel}
\email{yspennstate@gmail.com}
\date{November 2026}
\subjclass[2020]{60K37, 60G10, 82B20}
\keywords{random walk in random scenery, bad configurations, quasilocality, continuity of conditional probabilities}

\raggedbottom
\begin{document}
\begin{abstract}
Let $(S_t)_{t\in\Z}$ be a nearest-neighbour random walk on $\Z$ with $S_0=0$ and probability $p$ of a step to the
right, and let $(C_x)_{x\in\Z}$ be an independent field of independent fair signs. The reduced random walk in random
scenery is the colour process $Y_t=C_{S_t}$. Den Hollander and Steif conjectured that, under conditioning on the
colours at all nonzero times, every configuration is bad for $Y_0$ when $1/2\le p<4/5$ and every configuration is good
when $4/5<p\le1$. We prove that every configuration is good when $27/28<p\le1$. More precisely, the conditional
probabilities of $Y_0$ given the colours at the times $0<|t|\le m$ converge uniformly as $m\to\infty$, with an explicit
modulus, to a version of the conditional probability that is continuous everywhere. For colours with $\PP(C_x=1)=\alpha$
the same holds when $p>27/(27+4\beta^2)$, where $\beta=\min(\alpha,1-\alpha)$. The proof writes the posterior law of
the walk as a gas of disjoint time intervals separated by cut times, bounds the activities of the intervals, and swaps
the central part of two coupled samples between common cut times.
\end{abstract}
\maketitle

\section{Introduction}\label{sec:intro}

We describe the reduced random walk in random scenery, recall the notion of a bad configuration and the conjecture of
den Hollander and Steif, and state the result.

\subsection{The process}
Fix $p\in[1/2,1]$ and put $q=1-p$. Let $(X_t)_{t\in\Z}$ be independent with $\PP(X_t=1)=p$ and $\PP(X_t=-1)=q$, and let
$S_0=0$ and $S_t-S_{t-1}=X_t$ for every $t\in\Z$; thus $S_{-k}=-(X_{-k+1}+\dots+X_0)$ for $k\ge1$. Independently of
$(X_t)$, let $(C_x)_{x\in\Z}$ be independent with $\PP(C_x=1)=\alpha$ and $\PP(C_x=-1)=1-\alpha$, where
$\alpha\in(0,1)$, and write $\mu(1)=\alpha$, $\mu(-1)=1-\alpha$ and $\beta=\min(\alpha,1-\alpha)$. The \emph{reduced
random walk in random scenery} is the stationary process
\[
Y_t=C_{S_t},\qquad t\in\Z.
\]
Only $(Y_t)$ is observed, not the walk. The case $\alpha=1/2$ of fair colours is the one considered in
\cite{dHS06}. In ergodic theory random walk in random scenery appears through Kalikow's $T,T^{-1}$ transformation \cite{Kal82}; the
survey \cite{dHS06} describes its role as a source of non-Gibbsian measures.

\subsection{Bad configurations}
For $n\ge0$ let $\Lambda_n=\{-n,\dots,n\}$. For a configuration $\eta\in\{-1,1\}^{\Z\setminus\{0\}}$ and $m\ge1$ put
\begin{equation}\label{eq:gm}
g_m(\eta)=\PP\bigl(Y_0=1\mid Y_t=\eta_t\text{ for all }t\in\Lambda_m\setminus\{0\}\bigr).
\end{equation}
Every finite conditioning event has positive probability, since a monotone walk reads distinct sites. Following
\cite[Section~2.6]{dHS06} (see also \cite{MRM99,vEFS93}), $\eta$ is \emph{bad} for $Y_0$ if there is $\varepsilon>0$ such
that for every $n$ there are $m\ge n$ and a configuration $\delta$ equal to $\eta$ on $\Lambda_n\setminus\{0\}$ with
$|g_m(\eta)-g_m(\delta)|\ge\varepsilon$; otherwise $\eta$ is \emph{good}. The difference $|g_m(\eta)-g_m(\delta)|$ is the
total variation distance between the two conditional laws of $Y_0$. Every version of the conditional probability of
$Y_0$ given $(Y_t)_{t\ne0}$ is discontinuous at a bad configuration, and there is a version continuous at every good
configuration \cite{dHSvdW05,MRM99}.

\subsection{The conjecture and what is known}
Den Hollander and Steif \cite[Conjecture~2.16]{dHS06} conjectured that for fair colours every configuration is bad for
$Y_0$ when $1/2\le p<4/5$ and every configuration is good when $4/5<p\le1$. The theorems that precede the conjecture in
\cite[Section~2.6]{dHS06} concern the joint process of steps and colours observed outside time zero, and the earlier
paper \cite{dHSvdW05} classifies the bad configurations of that joint process; these results do not address the reduced
process. For the reduced process with \emph{one-sided} conditioning, on the colours at the times $t<0$ only, Blach\`ere,
den Hollander and Steif \cite[Theorem~1.3]{BdHS11} proved, for the walk that does not hold, that every configuration is bad for $1/2<p<4/5$, that good and bad
configurations coexist at $p=1/2$, and that every configuration is good for $p>1/(1+5^5/12^6)$, a value just above $0.99895$; walks
that hold are also treated there. The rest of the one-sided good half is \cite[Conjecture~1.4]{BdHS11}, and the phase
diagram of the reduced process is Question~B of Steif \cite{Ste21}.

\subsection{The result}

\begin{theorem}\label{thm:main}
Let $\alpha\in(0,1)$, $\beta=\min(\alpha,1-\alpha)$, and
\[
\frac{27}{27+4\beta^2}<p\le1 .
\]
Then $g_m$ converges uniformly on $\{-1,1\}^{\Z\setminus\{0\}}$ as $m\to\infty$ to a continuous function $g$, and $g$
is a version of $\PP(Y_0=1\mid(Y_t)_{t\ne0})$. Moreover there is a sequence $\varepsilon_n\to0$, depending only on
$p$ and $\alpha$ and given in \eqref{eq:modulus}, such that $|g_m(\eta)-g_m(\delta)|\le\varepsilon_n$ whenever $m\ge n$
and $\eta=\delta$ on $\Lambda_n\setminus\{0\}$. In particular every configuration is good for $Y_0$.
\end{theorem}

\begin{corollary}\label{cor:fair}
For fair colours, every configuration is good for $Y_0$ when $27/28<p\le1$.
\end{corollary}

Corollary~\ref{cor:fair} is the case $\beta=1/2$, and it proves the good half of \cite[Conjecture~2.16]{dHS06} for
$27/28<p\le1$. The range $4/5<p\le27/28$ and the bad half remain open.

\subsection{The value \texorpdfstring{$4/5$}{4/5}}
The value $4/5$ comes from a large-deviation balance which the proof below does not attempt to reach. A monochromatic
record of length $N$ can be read by a path with a fraction $v\ge1/2$ of right steps; such a path costs
$N(v\log(v/p)+(1-v)\log((1-v)/q))$ in the exponent and needs about $(2v-1)N$ new sites of the right colour, each at
cost $\log2$. The total $J_p(v)=v\log(v/p)+(1-v)\log((1-v)/q)+(2v-1)\log2$ is strictly convex with
$J_p'(1/2)=\log(4q/p)$, so its minimum on $[1/2,1]$ is at zero speed exactly when $p\le4/5$. Below $4/5$ a conditioned
walk can afford to return and reuse scenery, and distant colours can then influence the colour at the origin
\cite[Section~1.3]{BdHS11}. The proof of Theorem~\ref{thm:main} replaces this balance by a crude count, which costs the
gap between $4/5$ and $27/28$; see Remark~\ref{rem:sharp}.

\subsection{Idea of the proof}
Fix the conditioning window $\Lambda_m$. A time $k$ is a positive cut of the walk if the spatial edge
$\{S_k,S_k+1\}$ is crossed exactly once, from left to right, at time $k$. Splitting the window at the positive cuts
gives blocks with disjoint spatial ranges, so the posterior law of the walk, given any colours on
$\Lambda_m\setminus\{0\}$, is a hard-core gas of disjoint time intervals with local activities
(Lemma~\ref{lem:gas}). A block that contains no positive cut takes at least a third of its steps to the left
(Lemma~\ref{lem:third}), which bounds the activities by $\rho^L$ for an interval of $L$ steps, with $\rho<1$ exactly in
the range of Theorem~\ref{thm:main} (Lemma~\ref{lem:activity}). The gas is then dominated by independent intervals
(Lemma~\ref{lem:domination}), so two posterior walks, conditioned on configurations that agree on $\Lambda_n$, have a
common positive cut on each side of the origin inside $\Lambda_n$ with probability close to one
(Lemma~\ref{lem:common}). Exchanging the blocks between two common cuts preserves the joint weight and exchanges the
two colours at the origin, which proves the bound on $|g_m(\eta)-g_m(\delta)|$.

The full proof occupies Sections~\ref{sec:cuts}--\ref{sec:proof}. Section~\ref{sec:cuts} decomposes a path into blocks,
Section~\ref{sec:gas} writes the posterior law as an interval gas and bounds its activities, Section~\ref{sec:common}
finds common cuts, and Section~\ref{sec:proof} proves Theorem~\ref{thm:main}. Section~\ref{sec:remarks} discusses the
constant and the exact finite-window computations, and Appendix~\ref{app:notation} lists the notation.

\section{Positive cuts and blocks}\label{sec:cuts}

Fix integers $a<b$ and a nearest-neighbour path $s=(s_a,\dots,s_b)$. Call $k\in\{a,\dots,b-1\}$ a \emph{positive
cut} of $s$ if
\begin{equation}\label{eq:cut}
\max_{a\le j\le k}s_j\le s_k\qquad\text{and}\qquad\min_{k<j\le b}s_j>s_k .
\end{equation}
Then $s_{k+1}=s_k+1$, and $k$ is a positive cut exactly when the spatial edge $\{s_k,s_k+1\}$ is crossed only once by
$s$, at time $k$: a path that crosses this edge once, from left to right, stays at or to the left of $s_k$ before time
$k$ and at or to the right of $s_k+1$ after it, and conversely \eqref{eq:cut} rules out a second crossing.

Splitting $\{a,\dots,b\}$ at the positive cuts, that is, between $k$ and $k+1$ for every positive cut $k$, gives
consecutive \emph{blocks} $I=\{u,\dots,v\}$. Different blocks have disjoint spatial ranges, every block other than the
first starts at its minimum, and every block other than the last ends at its maximum.

\begin{lemma}\label{lem:blocks}
The map from a path to its sequence of blocks, each taken up to spatial translation, is a bijection onto the sequences
of consecutive blocks $I_1,\dots,I_J$ covering $\{a,\dots,b\}$, each carrying a nearest-neighbour path, such that
\begin{enumerate}
\item every block other than the first starts at its minimum and every block other than the last ends at its maximum;
\item no block has a positive cut of its own path.
\end{enumerate}
The path is recovered by joining consecutive blocks with one step to the right, and its positive cuts are exactly
these joining steps.
\end{lemma}

\begin{proof}
Let $s$ be a path. The blocks satisfy (1). A spatial edge inside the range of a block is crossed only by steps of that
block, since the ranges of different blocks are disjoint and the joining steps cross edges between ranges; hence a
positive cut of the path of a block would be a positive cut of $s$ inside the block, which does not exist. So (2)
holds. Conversely, given blocks with (1) and (2), translate each block so that it starts one site to the right of the
end of the previous block. By (1), the ranges are disjoint and ordered from left to right, and every joining step
satisfies \eqref{eq:cut}. An edge inside the range of a block is crossed only by that block, so by (2) no step inside
a block is a positive cut of the joined path. The two constructions are inverse to each other.
\end{proof}

The first block need not start at its minimum, the last need not end at its maximum, and a single block covering
$\{a,\dots,b\}$ has neither restriction; these cases are included in everything below. A block with one vertex is a
\emph{singleton}.

For a block $I=\{u,\dots,v\}$ write $L=v-u$ for its number of steps, $\ell$ for its number of steps to the left and
$d=s_v-s_u=L-2\ell$ for its displacement.

\begin{lemma}\label{lem:third}
For every block, $L\le3\ell$.
\end{lemma}

\begin{proof}
If $d\le0$ then $L=d+2\ell\le2\ell$. Let $d>0$. Each of the $d$ spatial edges between $s_u$ and $s_v$ is crossed by the
path of the block once more from left to right than from right to left. By Lemma~\ref{lem:blocks}(2) and the
characterisation of positive cuts after \eqref{eq:cut}, none of them is crossed only once, so each is crossed at least
once from right to left. These crossings are distinct left steps, hence $\ell\ge d$ and $L=d+2\ell\le3\ell$.
\end{proof}

\section{The interval gas}\label{sec:gas}

Fix $m\ge1$, put $a=-m$, $b=m$, and fix a configuration $\eta$ on $\Lambda_m\setminus\{0\}$. The law of
$(Y_t)_{t\in\Lambda_m}$ does not change if the walk is anchored at $S_a=0$ instead of $S_0=0$: translating the scenery
by $S_a$ gives, conditionally on the steps, a field with the same independent law. We use this anchoring throughout.
For a path $s$ on $\Lambda_m$ the probability of the steps is $p^{b-a-\ell(s)}q^{\ell(s)}$, where $\ell(s)$ is the number
of left steps, and the probability that $C_{s_t}=\eta_t$ for all $t\in\Lambda_m\setminus\{0\}$ is zero if two
observations at the same site disagree and otherwise $\prod_x\mu(\eta(x))$, the product over the distinct sites $x$
visited at an observed time, $\eta(x)$ being the common colour observed there. As reference we use the value for the
path that only steps to the right,
\[
B=p^{b-a}\prod_{t\in\Lambda_m\setminus\{0\}}\mu(\eta_t).
\]

For a block $I$ carrying a path $w$, let $h_I$ be the number of observed times in $I$, that is, of times in
$I\setminus\{0\}$, and let $r_I(w)$ be the number of distinct sites visited by $w$ at observed times. Define the
\emph{weight} of $(I,w)$ as zero if two observations in $I$ at the same site disagree, and otherwise
\begin{equation}\label{eq:weight}
A_I(w)=\Bigl(\frac qp\Bigr)^{\ell(w)}\prod_{t\in I\setminus\{0\}}\mu(\eta_t)^{-1}\prod_x\mu(\eta(x)),
\end{equation}
the last product over the $r_I(w)$ distinct sites visited at observed times in $I$. Let $w_I$ be the sum of $A_I(w)$
over the paths $w$ on $I$, up to translation, that satisfy the restrictions of Lemma~\ref{lem:blocks}: no positive cut,
and the boundary conditions in (1) that apply to $I$ according to whether $I$ contains $a$, $b$, both or neither. We call
$w_I$ the \emph{activity} of $I$. A singleton has activity $1$.

\begin{lemma}\label{lem:gas}
The probability of the event $\{Y_t=\eta_t,\ t\in\Lambda_m\setminus\{0\}\}$ equals
\begin{equation}\label{eq:gas}
B\sum_{\mathcal I}\prod_{I\in\mathcal I}w_I ,
\end{equation}
where $\mathcal I$ runs over the sets of pairwise disjoint intervals of $\Lambda_m$ with at least two vertices; the
vertices outside $\bigcup\mathcal I$ are singleton blocks. The same holds for the probability of the event
$\{Y_t=\eta_t,\ t\in\Lambda_m\setminus\{0\}\}\cap\{Y_0=1\}$ if, in the weights of the paths of the block that contains
$0$, time $0$ is treated as an observed time with colour $1$, so that the site $S_0$ joins the sites in the last product
of \eqref{eq:weight}, while the factor $\prod_{t\in I\setminus\{0\}}\mu(\eta_t)^{-1}$ is unchanged. For the singleton
block $\{0\}$ this activity is $\mu(1)$.
\end{lemma}

\begin{proof}
By Lemma~\ref{lem:blocks}, paths on $\Lambda_m$ correspond to sequences of blocks, and the sequence is determined by
the set of non-singleton blocks and the paths they carry. The joining steps are steps to the right, so the step
probability of the path, divided by $p^{b-a}$, is the product over blocks of $(q/p)^{\ell}$. The spatial ranges of
different blocks are disjoint, so the colour constraints of different blocks involve different sites, and the
probability of the colours factorises over blocks. Dividing by $B$, which is the product over blocks of
$\prod_{t\in I\setminus\{0\}}\mu(\eta_t)$, gives the product of the weights \eqref{eq:weight}. Summing over the paths of
each block gives \eqref{eq:gas}. For the second statement observe that the site $S_0$ belongs to the range of the block
containing $0$ only, so the colour at time $0$ is constrained only inside that block. If that site is also visited at
an observed time it must carry colour $1$, and its factor $\mu(1)$ already appears once in the product; if not, its
colour is independent of everything else and contributes the factor $\mu(1)$.
\end{proof}

Normalising \eqref{eq:gas} gives a probability measure on sets $\mathcal I$ of disjoint intervals, the \emph{posterior
interval gas}; its weights are the products of the activities, and there is no further interaction. Marking each
interval with its path, distributed proportionally to $A_I(w)$, and the block of $0$ with its colour at time $0$ gives
the posterior law of the walk and of $Y_0$. Lemma~\ref{lem:gas} is exact for every configuration $\eta$.

\begin{lemma}\label{lem:activity}
Let $\tau=q/(p\beta^2)$ and suppose $\tau<4/27$. Then every interval $I$ with $L\ge1$ steps has activity
\[
w_I\le\rho^L,\qquad \rho=\Bigl(\frac{27\tau}{4}\Bigr)^{1/3}<1 .
\]
The condition $\tau<4/27$ is equivalent to $p>27/(27+4\beta^2)$.
\end{lemma}

\begin{proof}
Let $w$ be a path on $I$ with weight $A_I(w)\ne0$, and let $R_I$ be the number of sites in its range. A
nearest-neighbour path visits every site of its range, and at most one vertex of $I$, the time $0$, is unobserved, so
\[
h_I=L+1-e_I,\qquad r_I(w)\ge R_I-e_I,\qquad R_I\ge|d|+1,
\]
where $e_I=1$ if $0\in I$ and $e_I=0$ otherwise. Hence $h_I-r_I(w)\le L+1-R_I\le L-|d|\le2\ell$. In
\eqref{eq:weight} the colour factors of the $r_I(w)$ distinct sites cancel against $r_I(w)$ of the $h_I$ factors
$\mu(\eta_t)^{-1}$, and each of the remaining $h_I-r_I(w)$ factors is at most $\beta^{-1}$. Therefore
\[
0\le A_I(w)\le\Bigl(\frac qp\Bigr)^{\ell}\beta^{-2\ell}=\tau^{\ell}.
\]
A path on $I$ is determined by its $L$ steps, and by Lemma~\ref{lem:third} it has $\ell\ge L/3$ left steps. Hence, for
every $z\ge1$,
\[
w_I\le\sum_{\ell\ge L/3}\binom L\ell\tau^{\ell}\le\sum_{\ell=0}^{L}\binom L\ell\tau^{\ell}z^{\ell-L/3}
=z^{-L/3}(1+\tau z)^L .
\]
Since $\tau<4/27<1/2$ we may take $z=1/(2\tau)$, which gives $w_I\le(2\tau)^{L/3}(3/2)^L=(27\tau/4)^{L/3}=\rho^L$.
Finally $27\tau<4$ means $27q<4p\beta^2$, that is, $p>27/(27+4\beta^2)$.
\end{proof}

For fair colours $\tau=4q/p$, $\rho^3=27q/p$, and the condition is $p>27/28$.

\begin{lemma}\label{lem:domination}
Let $\mathcal I$ be distributed according to a finite interval gas with activities $w_I\ge0$, and let $(\xi_I)$ be
independent Bernoulli variables with $\PP(\xi_I=1)=w_I/(1+w_I)$. Then $(\I\{I\in\mathcal I\})_I$ and $(\xi_I)_I$ can be
coupled so that $\I\{I\in\mathcal I\}\le\xi_I$ for every $I$.
\end{lemma}

\begin{proof}
Fix an interval $I$ and condition on the presence or absence of any set of other intervals. Removing $I$ maps every
admissible set containing $I$ injectively to an admissible set without $I$ that agrees on the conditioned intervals,
and divides its weight by $w_I$. Hence the conditional odds of $I\in\mathcal I$ are at most $w_I$, and its conditional
probability is at most $w_I/(1+w_I)$. Order the intervals, let $(U_I)$ be independent uniform variables on $[0,1]$, put
$\xi_I=\I\{U_I\le w_I/(1+w_I)\}$, and let $I\in\mathcal I$ if $U_I$ is at most the conditional probability of $I$ given
the earlier decisions. This gives the law of the gas and the required inequality.
\end{proof}

\section{Common cuts}\label{sec:common}

By Lemmas~\ref{lem:activity} and~\ref{lem:domination}, and since $x\mapsto x/(1+x)$ is increasing, the posterior
interval gas for any configuration is dominated by independent Bernoulli variables $\xi_I$ with
$\PP(\xi_I=1)=\rho^L/(1+\rho^L)$ for an interval of $L$ steps. A time edge $\{k,k+1\}$ of $\Lambda_m$ is \emph{covered}
by an interval that contains both $k$ and $k+1$; it is a positive cut exactly when no interval of $\mathcal I$ covers it.
For $1\le K<r$ put
\begin{equation}\label{eq:bound}
c=\exp\Bigl(-\frac{2\rho}{(1-\rho)^2}\Bigr),\qquad
\Phi(r,K)=(1-c)^{\lfloor r/(K+1)\rfloor}
+2\rho^{K+1}\Bigl(\frac{r+K+1}{1-\rho}+\frac{\rho}{(1-\rho)^2}\Bigr).
\end{equation}

\begin{lemma}\label{lem:common}
Assume $\rho<1$. Let $\mathcal I$ and $\mathcal I'$ be independent posterior interval gases on $\Lambda_m$ for two
configurations, and let $E$ be a set of $r$ consecutive time edges of $\Lambda_m$. For every $1\le K<r$, the probability
that no edge of $E$ is a positive cut of both walks is at most $\Phi(r,K)$.
\end{lemma}

\begin{proof}
Couple $\mathcal I$ and $\mathcal I'$ with independent dominating families $(\xi_I)$ and $(\xi'_I)$ as above. An edge
not covered by any interval with $\xi_I=1$ or $\xi'_I=1$ is a positive cut of both walks, so it suffices to bound the
probability that every edge of $E$ is covered in one of the two families. Choose $s=\lfloor r/(K+1)\rfloor$ edges of $E$
spaced $K+1$ apart. An interval of at most $K$ steps covers at most one of them. A fixed edge is covered by at most $L$
intervals of $L$ steps, so the probability that it is covered by no interval of at most $K$ steps in either family is
at least
\[
\prod_{L\ge1}\bigl(1+\rho^L\bigr)^{-2L}\ge\exp\Bigl(-2\sum_{L\ge1}L\rho^L\Bigr)=c ,
\]
using $\log(1+x)\le x$. These events for the $s$ chosen edges depend on disjoint sets of independent variables, so the
probability that each chosen edge is covered by a short interval is at most $(1-c)^s$. At most $r+L$ intervals of $L$
steps cover some edge of $E$, so the probability that an interval of more than $K$ steps with $\xi_I=1$ or $\xi'_I=1$
covers an edge of $E$ is at most
\[
2\sum_{L>K}(r+L)\rho^L=2\rho^{K+1}\Bigl(\frac{r+K+1}{1-\rho}+\frac{\rho}{(1-\rho)^2}\Bigr).
\]
If neither event occurs, some chosen edge is covered by no interval in either family. Intervals that do not fit inside
$\Lambda_m$ are absent, which only helps, so the bound does not depend on $m$ or on the configurations.
\end{proof}

With $K=\lfloor\sqrt r\rfloor$ both terms of $\Phi(r,K)$ tend to zero as $r\to\infty$, since $\lfloor r/(K+1)\rfloor\to\infty$
and $r\rho^{\sqrt r}\to0$. The constant $c$ can be very small; it affects the rate, not the convergence.

\section{Proof of Theorem~\ref{thm:main}}\label{sec:proof}

For $p=1$ the walk is $S_t=t$, so $(Y_t)$ is an independent sequence and $g_m\equiv\alpha$. Assume
$27/(27+4\beta^2)<p<1$, so that $\rho<1$ by Lemma~\ref{lem:activity}. Put
\begin{equation}\label{eq:modulus}
\varepsilon_n=\min\bigl(1,\,2\Phi(n,\lfloor\sqrt n\rfloor)\bigr)\quad(n\ge2),\qquad\varepsilon_0=\varepsilon_1=1 .
\end{equation}

Let $m\ge n\ge2$ and let $\eta$, $\delta$ be configurations that agree on $\Lambda_n\setminus\{0\}$. Sample
independently the marked posterior interval gases on $\Lambda_m$ for $\eta$ and for $\delta$, as after
Lemma~\ref{lem:gas}, and let $U$ and $V$ be the two colours at time $0$, so that $\PP(U=1)=g_m(\eta)$ and
$\PP(V=1)=g_m(\delta)$. Let $G$ be the event that the two walks have a common positive cut among the $n$ edges
$\{k,k+1\}$ with $-n\le k\le-1$, and another among the $n$ edges with $0\le k\le n-1$. On $G$ let $k_-$ be the largest
common positive cut with $k_-\le-1$ and $k_+$ the smallest with $k_+\ge0$, and exchange between the two samples all the
marked blocks contained in $\{k_-+1,\dots,k_+\}$, including the colours at time $0$.

These blocks lie in $\Lambda_n$, where the two configurations agree, and none of them is the first or the last block,
so by \eqref{eq:weight} each has the same weight in both samples, and the joint weight of the pair is preserved. The
exchange does not change the set of common positive cuts, since inside $\{k_-+1,\dots,k_+\}$ it exchanges the two cut
sets and leaves their intersection unchanged. So it maps $G$ to itself, chooses the same $k_\pm$ when applied again, and
is a weight-preserving involution of $G$ that exchanges $U$ and $V$. Therefore
$\E[(\I\{U=1\}-\I\{V=1\})\I_G]=0$, and by Lemma~\ref{lem:common} applied to each side,
\[
|g_m(\eta)-g_m(\delta)|=\bigl|\E[(\I\{U=1\}-\I\{V=1\})\I_{G^c}]\bigr|\le\PP(G^c)\le\varepsilon_n .
\]
This is the bound of Theorem~\ref{thm:main}, uniform in $m\ge n$ and in the two configurations; for $n\le1$ it holds
trivially. Since $\varepsilon_n\to0$, no configuration is bad.

For $m\ge n$, $g_n(\eta)$ is an average of the values $g_m(\delta)$ over configurations $\delta$ that agree with $\eta$
on $\Lambda_n\setminus\{0\}$, so $|g_m(\eta)-g_n(\eta)|\le\varepsilon_n$. Hence $(g_m)$ is uniformly Cauchy and converges
uniformly to a function $g$, which is continuous because each $g_m$ depends on finitely many coordinates. By the
martingale convergence theorem, $g_m((Y_t)_{t\ne0})\to\PP(Y_0=1\mid(Y_t)_{t\ne0})$ almost surely, so $g$ is a version of
this conditional probability. \qed

\section{Remarks}\label{sec:remarks}

We comment on the constant and on the finite-window computations that accompany the paper.

\begin{remark}\label{rem:sharp}
The constant $27/28$ comes from two losses. Lemma~\ref{lem:third} uses only that a block without a positive cut spends
a third of its steps going left, and Lemma~\ref{lem:activity} then counts all step sequences with that many left steps,
ignoring that most of them do have positive cuts. The colour bound $\beta^{-1}$ per repeated observation is also
attained only by records that match the walk perfectly. Reaching $4/5$ would need control of the activities of
conditioned blocks close to the large-deviation balance described in the introduction, or a different separation
argument. Theorem~\ref{thm:main} says nothing about the bad half of the conjecture. The coexistence of good and bad
configurations at $p=1/2$ in the one-sided setting \cite[Theorem~1.3(iii)]{BdHS11} is not a counterexample to the
two-sided conjecture, since observations at positive times can change the argument.
\end{remark}

\begin{remark}\label{rem:general}
For general colours the large-deviation heuristic of the introduction, applied to a record of the less likely colour,
with cost $\log(1/\beta)$ per new site, suggests the boundary $p=1/(1+\beta^2)$ in place of $4/5$; for $\beta=1/2$ this is $4/5$. The bound
$27/(27+4\beta^2)$ of Theorem~\ref{thm:main} lies above it for every $\beta$.
\end{remark}

\begin{remark}\label{rem:finite}
For rational $p$ and fair colours the probabilities $g_m(\eta)$ can be computed exactly. Anchor the walk at time $-m$.
Since the walk has nearest-neighbour steps, its range is an interval that grows at an endpoint, so a state consisting of the
position within the range, the length of the range and the colours on it, up to translation, carries all the
information needed; each step multiplies the weight by $p$ or $q$, a revisit with the observed colour by $1$, and a
new site by $1/2$, and a revisit with the other colour is discarded. Running this transfer with both colours at time $0$
gives $g_m(\eta)=Z_+/(Z_++Z_-)$ in exact rational arithmetic. We computed $|g_m(\eta)-g_m(\delta)|$ for $\eta$ equal to
$1$ at every nonzero time and four choices of $\delta$ agreeing with $\eta$ on $\Lambda_n\setminus\{0\}$, $n=2,4,6$,
$m-n\in\{0,2,4,8,12,20,28\}$, and ten values of $p$ from $1/2$ to $199/200$: 840 comparisons with $m\le34$. For
$\delta$ repeating the pattern $-1,-1,1,1$ outwards on both sides outside $\Lambda_n$ and $m=n+28$, the differences are
the following.
\begin{center}
\begin{tabular}{lccc}
\toprule
$p$ & $n=2$ & $n=4$ & $n=6$\\
\midrule
$1/2$ & $0.2399$ & $0.0867$ & $0.0357$\\
$4/5$ & $0.2321$ & $0.1082$ & $0.0595$\\
$17/20$ & $0.2475$ & $0.1226$ & $0.0721$\\
$97/100$ & $0.0668$ & $0.0141$ & $0.0035$\\
$199/200$ & $0.0102$ & $0.0004$ & $0.0000$\\
\bottomrule
\end{tabular}
\end{center}
At $p=4/5$ and $n=4$ the difference grows with $m$, from $0.045$ at $m=6$ to $0.108$ at $m=32$. Badness concerns
arbitrarily large $n$ and an optimisation over all alterations, so these numbers neither prove nor refute anything
below $27/28$; they show only that the windows reachable by exact computation are far from the asymptotic regime. The
code also checks the transfer against direct enumeration of all step sequences and against the interval-gas formula
of Lemma~\ref{lem:gas} (417 exact checks), and it checks the block structure of Lemma~\ref{lem:blocks},
Lemma~\ref{lem:third} and the bound $h_I-r_I(w)\le2\ell$ on the $58{,}474$ blocks of all $32{,}767$ step sequences of
length at most $14$.
\end{remark}

\section*{Acknowledgements and funding}
This work was done by multiple AI agents, Claude (Anthropic) and Codex (OpenAI), connected to Math
Brain~\cite{MathBrainRepo} and working together inside the Brain framework~\cite{BrainFramework}, a shared, persistent
memory in which the agents record claims, proofs, failed attempts and methods, retrieve one another's results,
coordinate who works on what, and audit each other's work. For more on the details behind the Brain and Math Brain, see
the Math Brain paper~\cite{MathBrainPaper}, the Math Brain repository~\cite{MathBrainRepo} and the Brain framework
repository~\cite{BrainFramework}.

The research presented in this paper was supported by the European Research Council (ERC) under the European Union's
Horizon Europe research and innovation programme (grant agreement No.~101041711), by the Simons Foundation as part of
the Collaboration on the Mathematical and Scientific Foundations of Deep Learning, by Heights Labs, by Convex Nexus
Capital, by the Israel Science Foundation (grant number 2258/19), and by the Israel Science Foundation (ISF Grant
4101/25). Computational resources were provided in part by Google Cloud through its education credits programme (credit
530550103).

\section*{Declaration of competing interest}
The author holds equity in Heights Labs and in Convex Nexus Capital and has a management role at Convex Nexus Capital.
These entities are included in the funding declaration above.

\section*{Declaration of generative AI and AI-assisted technologies}
A Codex agent (OpenAI) proved Theorem~\ref{thm:main} and wrote the code of Remark~\ref{rem:finite}; a second Codex agent
checked the proof and the literature as a referee; Claude agents (Anthropic) checked the steps and prepared the
manuscript. The author directed the work and is responsible for the content of the paper and for decisions concerning
its dissemination.

\section*{Code availability}
The code and data of Remark~\ref{rem:finite} are in the repository
\begin{center}\url{https://github.com/yspennstate/rwrs-two-sided-good-configurations}.\end{center}
They use Python~3 and its standard library only; all probabilities are stored as exact fractions.

\appendix
\section{Notation}\label{app:notation}

\begin{center}
\begin{tabular}{p{3.0cm}p{8.0cm}p{3.1cm}}
\toprule
Symbol & Meaning & First appears\\
\midrule
$p$, $q$ & probabilities of a step to the right and to the left & Section~\ref{sec:intro}\\
$\PP$, $\E$ & probability and expectation & Section~\ref{sec:intro}\\
$X_t$, $S_t$ & steps and position of the walk & Section~\ref{sec:intro}\\
$C_x$, $\alpha$, $\mu$ & scenery, its probability of colour $1$, its law & Section~\ref{sec:intro}\\
$\beta$ & $\min(\alpha,1-\alpha)$ & Section~\ref{sec:intro}\\
$Y_t$ & the reduced process $C_{S_t}$ & Section~\ref{sec:intro}\\
$\Lambda_n$ & $\{-n,\dots,n\}$ & \eqref{eq:gm}\\
$\eta$, $\delta$ & configurations on $\Z\setminus\{0\}$ & \eqref{eq:gm}\\
$g_m$, $g$ & conditional probability of $Y_0=1$ given $\Lambda_m\setminus\{0\}$, its limit & \eqref{eq:gm},
Theorem~\ref{thm:main}\\
$\varepsilon_n$ & the modulus of Theorem~\ref{thm:main} & \eqref{eq:modulus}\\
$N$, $v$, $J_p$ & record length, fraction of right steps, large-deviation cost & Section~\ref{sec:intro}\\
$\I$ & indicator & Lemma~\ref{lem:domination}\\
$a$, $b$ & endpoints of the time window & Section~\ref{sec:cuts}\\
$I=\{u,\dots,v\}$ & a block & Section~\ref{sec:cuts}\\
$L$, $\ell$, $d$ & steps, left steps and displacement of a block & Section~\ref{sec:cuts}\\
$B$ & weight of the path that only steps right & Section~\ref{sec:gas}\\
$\eta(x)$ & colour observed at the site $x$ & Section~\ref{sec:gas}\\
$h_I$, $r_I(w)$ & observed times in $I$, distinct sites visited at them & Section~\ref{sec:gas}\\
$A_I(w)$, $w_I$ & weight of a block path, activity of an interval & \eqref{eq:weight}\\
$\mathcal I$ & set of non-singleton blocks & \eqref{eq:gas}\\
$R_I$, $e_I$ & size of the range, indicator that $0\in I$ & Lemma~\ref{lem:activity}\\
$\tau$, $\rho$ & $q/(p\beta^2)$, $(27\tau/4)^{1/3}$ & Lemma~\ref{lem:activity}\\
$z$ & auxiliary parameter of the binomial bound & Lemma~\ref{lem:activity}\\
$\xi_I$, $U_I$ & dominating Bernoulli and uniform variables & Lemma~\ref{lem:domination}\\
$c$, $\Phi(r,K)$ & constants of the common-cut bound & \eqref{eq:bound}\\
$E$, $r$, $K$, $s$ & a set of $r$ edges, a length cut-off, number of chosen edges & Lemma~\ref{lem:common}\\
$U$, $V$, $G$, $k_\pm$ & colours at time $0$, the event of common cuts, the chosen cuts & Section~\ref{sec:proof}\\
$Z_\pm$ & weights with colour $\pm1$ at time $0$ & Remark~\ref{rem:finite}\\
\bottomrule
\end{tabular}
\end{center}

\begin{thebibliography}{99}
\bibitem{BdHS11} S.~Blach\`ere, F.~den Hollander and J.~E.~Steif, A crossover for the bad configurations of random walk
in random scenery, \emph{Ann.\ Probab.} \textbf{39} (2011), no.~5, 2018--2041. \doi{10.1214/11-AOP664}
\bibitem{dHS06} F.~den Hollander and J.~E.~Steif, Random walk in random scenery: a survey of some recent results, in
\emph{Dynamics \& Stochastics}, IMS Lecture Notes Monogr.\ Ser.\ \textbf{48} (2006), 53--65.
\doi{10.1214/074921706000000077}
\bibitem{dHSvdW05} F.~den Hollander, J.~E.~Steif and P.~van der Wal, Bad configurations for random walk in random scenery
and related subshifts, \emph{Stochastic Process.\ Appl.} \textbf{115} (2005), no.~7, 1209--1232.
\doi{10.1016/j.spa.2005.03.001}
\bibitem{Kal82} S.~A.~Kalikow, $T,T^{-1}$ transformation is not loosely Bernoulli, \emph{Ann.\ of Math.\ (2)} \textbf{115}
(1982), no.~2, 393--409. \doi{10.2307/1971397}
\bibitem{MRM99} C.~Maes, F.~Redig and A.~Van Moffaert, Almost Gibbsian versus weakly Gibbsian measures,
\emph{Stochastic Process.\ Appl.} \textbf{79} (1999), no.~1, 1--15. \doi{10.1016/S0304-4149(98)00083-0}
\bibitem{Ste21} J.~E.~Steif, Some results and open questions concerning Gibbs states, quasilocality and bilateral
determinism, in Mini-Workshop: One-sided and Two-sided Stochastic Descriptions, \emph{Oberwolfach Rep.} \textbf{17}
(2020), no.~1, 632--633. \doi{10.4171/OWR/2020/11}
\bibitem{vEFS93} A.~C.~D.~van Enter, R.~Fern\'andez and A.~D.~Sokal, Regularity properties and pathologies of
position-space renormalization-group transformations: scope and limitations of Gibbsian theory, \emph{J.\ Stat.\ Phys.}
\textbf{72} (1993), no.~5--6, 879--1167. \doi{10.1007/BF01048183}
\bibitem{MathBrainRepo} Y.~Shmalo, Math Brain: public snapshot, code and tests, 2026,
\url{https://github.com/yspennstate/math-brain-public}.
\bibitem{MathBrainPaper} Y.~Shmalo, \emph{Math Brain: shared, auditable memory for multi-agent mathematical research},
October 2026, \url{https://github.com/yspennstate/math-brain-public/tree/main/papers/math-brain-2026-10-24}.
\bibitem{BrainFramework} Y.~Shmalo, Brain Framework, public software repository, 2026,
\url{https://github.com/yspennstate/brain-framework-public}.
\end{thebibliography}

\end{document}
